\begin{abstract}
Outcome-only reinforcement learning with verifiable rewards (RLVR) improves final-answer accuracy but provides weak supervision for intermediate reasoning quality.
Learned process reward models can supply denser guidance, yet they require costly annotations or introduce additional model uncertainty.
We present \textbf{TopoPRM}, a two-part framework that combines deterministic process supervision with efficient reasoning transfer.

First, we introduce a \emph{Verifiable Process Reward Model} that maps free-form reasoning traces to inferred dependency DAGs via a rule-based extractor.
On these graphs we compute two complementary deterministic rewards: a \emph{topology reward} for global structural regularity (acyclicity, dependency consistency) and a \emph{continuity reward} for local step-to-step traceability.
Both rewards are verifiable in an operational sense---reproducible, auditable, and programmatically checkable---without claiming complete semantic proof verification.
A correctness-first stratified shaping mechanism ensures that answer correctness remains the primary optimisation target.

Second, we introduce \emph{reverse-KL reasoning distillation}: a teacher trained with TopoPRM rewards is filtered by process-aware quality criteria, and a smaller student is trained via reverse KL to concentrate on high-quality reasoning modes, internalising graph-enriched behaviour as compact chain-like reasoning.

On Chinese mathematical critique benchmarks, TopoPRM improves overall critique accuracy by $+13.4$ points over outcome-only GRPO while reducing mean response length by 11\%.
The distilled Qwen3-8B student achieves 82.5\% on GSM8K and 59.8\% on MATH-500, competitive with models of comparable scale trained on substantially more data.
\end{abstract}
